\documentclass[../main.tex]{subfiles}

\begin{document}
\section{Introduction}

\subsection{Aim}
\noindent
This report compares three \gls{ML} algorithms on a medical image
classification task using the PathMNIST subset of the MedMNIST v2
benchmark~\cite{yang2023medmnistv2}. PathMNIST consists of 28$\times$28
\gls{HE}-stained colour images of nine colorectal tissue types and is well
suited to a controlled comparison between a classical kernel method, a fully
connected feed-forward network and a convolutional architecture.

\noindent
Two of the three algorithms are prescribed by the assignment specification:
the \gls{MLP} and the \gls{CNN}. From the first six weeks of the course we
chose the \gls{SVM} as the third method, as it offers a non-linear yet
mathematically transparent decision boundary that contrasts well with the
representation learning performed by neural networks.

\noindent
The aim of the study is therefore three-fold: (i) to design and implement
each algorithm at a scale that fits the available hardware while preserving
the inductive biases that distinguish the three model families; (ii) to
perform a disciplined hyper-parameter search for each model on a held-out
validation split, recording both performance and runtime so the trade-offs
are visible; and (iii) to evaluate the best configuration of each model on
the PathMNIST test set and use the differences in accuracy, macro-$F_{1}$,
runtime and confusion patterns to discuss why each model behaves the way it
does.

\subsection{Importance}
\noindent
Histopathological assessment of \gls{HE}-stained slides is a corner-stone of
oncology diagnosis, but manual scoring is labour-intensive and shows
substantial inter-observer variability. Computational pathology aims to
automate parts of this pipeline; the resulting models can serve as
second-opinion tools, accelerate annotation for downstream studies, and
surface morphological features that correlate with prognosis or treatment
response.

\noindent
Comparing classifiers on this task matters beyond raw accuracy. A clinically
useful tool also needs to be \emph{predictable} (its failure modes should be
intelligible to a pathologist), \emph{efficient} (so it can run on commodity
hardware in a hospital setting), and \emph{calibrated} to per-class
performance, because rare tissue types may carry the highest diagnostic
weight. The three classifiers in this study sit at very different points on
this trade-off space: the \gls{SVM} is interpretable through its support
vectors and decision margin but scales poorly with sample size; the \gls{MLP}
is fully differentiable end-to-end but ignores spatial structure; the
\gls{CNN} exploits translation invariance and locality but is more opaque.
Quantifying these differences on a real medical dataset, even a small one
such as PathMNIST, gives concrete guidance for the design of larger systems
and is the contribution of this report.

\end{document}
